\documentclass[aspectratio=169]{beamer}
\usepackage[utf8]{inputenc}
\usepackage[russian]{babel}
\usepackage{graphicx}
\usepackage{amsmath}
\usepackage{multirow}
\usepackage{xcolor}

% Выбор темы
\usetheme{Madrid}
\usecolortheme{default}

% Информация о презентации
\title[Таблицы в LaTeX]{Создание и форматирование таблиц в LaTeX}
\subtitle{Лабораторная работа №5}
\author[Викторов Е.И.]{Викторов Егор Игоревич\\Группа НПМмд02-24}
\institute[РУДН]{Федеральное государственное образовательное учреждение\\высшего образования\\«Российский Университет Дружбы Народов имени Патриса Лумумбы»}
\date{\today}

\begin{document}

% Титульный слайд
\begin{frame}
    \titlepage
\end{frame}

% Содержание
\begin{frame}{Содержание}
    \tableofcontents
\end{frame}

% Раздел 1: Введение
\section{Введение}

\begin{frame}{Цель работы}
    \begin{block}{Основная цель}
        Изучить основные принципы создания и форматирования таблиц в системе компьютерной вёрстки LaTeX
    \end{block}
    
    \vspace{0.5cm}
    
    \begin{itemize}
        \item Освоить окружение \texttt{tabular}
        \item Изучить команды выравнивания содержимого ячеек
        \item Освоить объединение столбцов с помощью \texttt{\textbackslash multicolumn}
        \item Научиться создавать профессионально оформленные таблицы
    \end{itemize}
\end{frame}

\begin{frame}{Задачи работы}
    \begin{enumerate}
        \item Создать простые таблицы с различными типами выравнивания
        \item Исследовать поведение LaTeX при несоответствии количества элементов
        \item Освоить команду \texttt{\textbackslash multicolumn}
        \item Создать сложные таблицы с комбинированным форматированием
        \item Проанализировать результаты и сформулировать выводы
    \end{enumerate}
\end{frame}

% Раздел 2: Теоретическая часть
\section{Теоретическая часть}

\begin{frame}[fragile]{Основы создания таблиц}
    \begin{block}{Базовый синтаксис}
        \begin{verbatim}
\begin{tabular}{спецификация_столбцов}
    содержимое таблицы
\end{tabular}
        \end{verbatim}
    \end{block}
    
    \vspace{0.5cm}
    
    \begin{exampleblock}{Пример}
        \begin{center}
        \begin{tabular}{|l|c|r|}
        \hline
        Левый & Центр & Правый \\
        \hline
        Текст & Текст & Текст \\
        \hline
        \end{tabular}
        \end{center}
    \end{exampleblock}
\end{frame}

\begin{frame}{Спецификация столбцов}
    \begin{columns}[T]
        \begin{column}{0.5\textwidth}
            \begin{block}{Типы выравнивания}
                \begin{itemize}
                    \item \texttt{l} — по левому краю
                    \item \texttt{c} — по центру
                    \item \texttt{r} — по правому краю
                    \item \texttt{|} — вертикальная линия
                \end{itemize}
            \end{block}
        \end{column}
        
        \begin{column}{0.5\textwidth}
            \begin{block}{Разделители}
                \begin{itemize}
                    \item \texttt{\&} — между ячейками
                    \item \texttt{\textbackslash\textbackslash} — конец строки
                    \item \texttt{\textbackslash hline} — горизонтальная линия
                \end{itemize}
            \end{block}
        \end{column}
    \end{columns}
    
    \vspace{0.5cm}
    
    \begin{alertblock}{Важно!}
        Количество элементов в строке должно соответствовать количеству объявленных столбцов
    \end{alertblock}
\end{frame}

\begin{frame}[fragile]{Команда multicolumn}
    \begin{block}{Синтаксис}
        \begin{verbatim}
\multicolumn{количество}{выравнивание}{содержимое}
        \end{verbatim}
    \end{block}
    
    \vspace{0.3cm}
    
    \begin{exampleblock}{Пример использования}
        \begin{center}
        \begin{tabular}{|l|c|r|}
        \hline
        \multicolumn{3}{|c|}{\textbf{Заголовок}} \\
        \hline
        Имя & Возраст & Город \\
        \hline
        Иван & 25 & Москва \\
        \hline
        \end{tabular}
        \end{center}
    \end{exampleblock}
\end{frame}

% Раздел 3: Практическая часть
\section{Практическая часть}

\begin{frame}{Эксперимент 1: Выравнивание по левому краю}
    \begin{columns}[T]
        \begin{column}{0.6\textwidth}
            \begin{center}
            \begin{tabular}{|l|l|l|}
            \hline
            \textbf{Имя} & \textbf{Фамилия} & \textbf{Возраст} \\
            \hline
            Иван & Иванов & 25 \\
            Мария & Петрова & 30 \\
            Александр & Сидоров & 28 \\
            \hline
            \end{tabular}
            \end{center}
        \end{column}
        
        \begin{column}{0.4\textwidth}
            \begin{block}{Результат}
                Все данные выровнены по левому краю
            \end{block}
            
            \begin{alertblock}{Применение}
                Текстовая информация, названия
            \end{alertblock}
        \end{column}
    \end{columns}
\end{frame}

\begin{frame}{Эксперимент 2: Выравнивание по центру}
    \begin{columns}[T]
        \begin{column}{0.6\textwidth}
            \begin{center}
            \begin{tabular}{|c|c|c|}
            \hline
            \textbf{Товар} & \textbf{Количество} & \textbf{Цена} \\
            \hline
            Яблоки & 10 & 100 \\
            Груши & 5 & 150 \\
            Апельсины & 8 & 200 \\
            \hline
            \end{tabular}
            \end{center}
        \end{column}
        
        \begin{column}{0.4\textwidth}
            \begin{block}{Результат}
                Содержимое по центру ячеек
            \end{block}
            
            \begin{alertblock}{Применение}
                Заголовки, симметричное отображение
            \end{alertblock}
        \end{column}
    \end{columns}
\end{frame}

\begin{frame}{Эксперимент 3: Выравнивание по правому краю}
    \begin{columns}[T]
        \begin{column}{0.6\textwidth}
            \begin{center}
            \begin{tabular}{|r|r|r|}
            \hline
            \textbf{Год} & \textbf{Доход} & \textbf{Расход} \\
            \hline
            2021 & 1000000 & 800000 \\
            2022 & 1200000 & 900000 \\
            2023 & 1500000 & 1000000 \\
            \hline
            \end{tabular}
            \end{center}
        \end{column}
        
        \begin{column}{0.4\textwidth}
            \begin{block}{Результат}
                Числа выровнены по правому краю
            \end{block}
            
            \begin{alertblock}{Применение}
                Числовые данные, финансовые показатели
            \end{alertblock}
        \end{column}
    \end{columns}
\end{frame}

\begin{frame}{Эксперимент 4: Комбинированное выравнивание}
    \begin{center}
    \begin{tabular}{|l|c|r|}
    \hline
    \textbf{Название (слева)} & \textbf{Категория (центр)} & \textbf{Цена (справа)} \\
    \hline
    Ноутбук & Электроника & 50000 \\
    Книга & Литература & 500 \\
    Стул & Мебель & 3000 \\
    \hline
    \end{tabular}
    \end{center}
    
    \vspace{0.5cm}
    
    \begin{block}{Вывод}
        Каждый столбец имеет оптимальное выравнивание для своего типа данных
    \end{block}
\end{frame}

\begin{frame}{Эксперимент 5: Недостаточное количество элементов}
    \begin{center}
    \begin{tabular}{|l|c|r|r|}
    \hline
    \textbf{Столбец 1} & \textbf{Столбец 2} & \textbf{Столбец 3} & \textbf{Столбец 4} \\
    \hline
    Полная строка & Данные & 100 & 200 \\
    \hline
    Только два & элемента & & \\
    \hline
    Один & & & \\
    \hline
    \end{tabular}
    \end{center}
    
    \vspace{0.5cm}
    
    \begin{exampleblock}{Результат}
        \begin{itemize}
            \item LaTeX успешно компилирует таблицу
            \item Пустые ячейки остаются справа
            \item Ошибок компиляции не возникает
        \end{itemize}
    \end{exampleblock}
\end{frame}

\begin{frame}{Эксперимент 6: Избыточное количество элементов}
    \begin{alertblock}{Ошибка компиляции}
        При попытке указать больше элементов, чем объявлено столбцов:
        
        \vspace{0.3cm}
        
        \texttt{! Extra alignment tab has been changed to \textbackslash cr.}
    \end{alertblock}
    
    \vspace{0.5cm}
    
    \begin{block}{Причина}
        Количество элементов в строке превышает количество объявленных столбцов
    \end{block}
    
    \vspace{0.3cm}
    
    \begin{exampleblock}{Решение}
        Строго соблюдать соответствие между объявленным и фактическим количеством столбцов
    \end{exampleblock}
\end{frame}

\begin{frame}{Эксперимент 7: Базовое использование multicolumn}
    \begin{center}
    \begin{tabular}{|l|c|r|}
    \hline
    \multicolumn{3}{|c|}{\textbf{Заголовок таблицы}} \\
    \hline
    \textbf{Имя} & \textbf{Возраст} & \textbf{Город} \\
    \hline
    Иван & 25 & Москва \\
    Мария & 30 & Санкт-Петербург \\
    \hline
    \end{tabular}
    \end{center}
    
    \vspace{0.5cm}
    
    \begin{block}{Результат}
        Три столбца объединены в один с центральным выравниванием для создания общего заголовка
    \end{block}
\end{frame}

\begin{frame}{Эксперимент 8: Объединение двух столбцов}
    \begin{center}
    \begin{tabular}{|l|c|c|r|}
    \hline
    \textbf{Продукт} & \multicolumn{2}{c|}{\textbf{Характеристики}} & \textbf{Цена} \\
    \hline
    Телефон & Память & 128 ГБ & 30000 \\
    Планшет & Экран & 10 дюймов & 25000 \\
    \hline
    \end{tabular}
    \end{center}
    
    \vspace{0.5cm}
    
    \begin{block}{Результат}
        Два средних столбца объединены для создания общего заголовка «Характеристики»
    \end{block}
\end{frame}

\begin{frame}{Эксперимент 9: Сложная таблица}
    \begin{center}
    \scalebox{0.85}{
    \begin{tabular}{|l|c|c|c|c|}
    \hline
    \multicolumn{5}{|c|}{\textbf{Отчёт за 2023 год}} \\
    \hline
    \multicolumn{1}{|c|}{\textbf{Квартал}} & \multicolumn{2}{c|}{\textbf{Доходы}} & \multicolumn{2}{c|}{\textbf{Расходы}} \\
    \hline
    & \textbf{План} & \textbf{Факт} & \textbf{План} & \textbf{Факт} \\
    \hline
    Q1 & 100 & 110 & 80 & 85 \\
    Q2 & 120 & 115 & 90 & 88 \\
    Q3 & 130 & 140 & 95 & 100 \\
    Q4 & 150 & 155 & 100 & 105 \\
    \hline
    \multicolumn{1}{|c|}{\textbf{Итого:}} & \multicolumn{2}{c|}{\textbf{520}} & \multicolumn{2}{c|}{\textbf{378}} \\
    \hline
    \end{tabular}
    }
    \end{center}
    
    \vspace{0.3cm}
    
    \begin{block}{Результат}
        Многоуровневая таблица с объединёнными ячейками на разных уровнях
    \end{block}
\end{frame}

% Раздел 4: Результаты
\section{Результаты исследования}

\begin{frame}{Сравнительная таблица типов выравнивания}
    \begin{center}
    \begin{tabular}{|c|c|l|}
    \hline
    \textbf{Тип} & \textbf{Обозначение} & \textbf{Оптимальное применение} \\
    \hline
    Левое & \texttt{l} & Текстовые данные, названия \\
    \hline
    Центральное & \texttt{c} & Заголовки, короткие значения \\
    \hline
    Правое & \texttt{r} & Числовые данные, суммы \\
    \hline
    \end{tabular}
    \end{center}
    
    \vspace{0.5cm}
    
    \begin{block}{Рекомендация}
        Выбирайте тип выравнивания в зависимости от характера данных в столбце
    \end{block}
\end{frame}

\begin{frame}{Анализ ошибок}
    \begin{center}
    \begin{tabular}{|l|l|c|}
    \hline
    \textbf{Ситуация} & \textbf{Результат} & \textbf{Ошибка?} \\
    \hline
    Мало элементов & Пустые ячейки & \textcolor{green}{Нет} \\
    \hline
    Много элементов & Не компилируется & \textcolor{red}{Да} \\
    \hline
    Правильное количество & Нормальная таблица & \textcolor{green}{Нет} \\
    \hline
    \end{tabular}
    \end{center}
    
    \vspace{0.5cm}
    
    \begin{alertblock}{Важный вывод}
        Всегда проверяйте соответствие количества ячеек и объявленных столбцов!
    \end{alertblock}
\end{frame}

\begin{frame}{Преимущества использования multicolumn}
    \begin{columns}[T]
        \begin{column}{0.5\textwidth}
            \begin{block}{Возможности}
                \begin{itemize}
                    \item Объединение столбцов
                    \item Создание заголовков
                    \item Многоуровневая структура
                    \item Гибкое форматирование
                \end{itemize}
            \end{block}
        \end{column}
        
        \begin{column}{0.5\textwidth}
            \begin{exampleblock}{Применение}
                \begin{itemize}
                    \item Сложные отчёты
                    \item Научные таблицы
                    \item Финансовые документы
                    \item Статистические данные
                \end{itemize}
            \end{exampleblock}
        \end{column}
    \end{columns}
\end{frame}

% Раздел 5: Выводы
\section{Выводы}

\begin{frame}{Основные выводы (1/2)}
    \begin{enumerate}
        \item \textbf{Типы выравнивания:}
        \begin{itemize}
            \item Левое (\texttt{l}) — для текста
            \item Центральное (\texttt{c}) — для заголовков
            \item Правое (\texttt{r}) — для чисел
            \item Комбинирование повышает читаемость
        \end{itemize}
        
        \vspace{0.3cm}
        
        \item \textbf{Обработка ошибок:}
        \begin{itemize}
            \item Недостаток элементов — пустые ячейки (без ошибок)
            \item Избыток элементов — ошибка компиляции
            \item Важно соблюдать соответствие количества столбцов
        \end{itemize}
    \end{enumerate}
\end{frame}

\begin{frame}{Основные выводы (2/2)}
    \begin{enumerate}
        \setcounter{enumi}{2}
        \item \textbf{Команда multicolumn:}
        \begin{itemize}
            \item Гибкое объединение столбцов
            \item Создание многоуровневых заголовков
            \item Индивидуальное выравнивание для каждого объединения
        \end{itemize}
        
        \vspace{0.3cm}
        
        \item \textbf{Практические навыки:}
        \begin{itemize}
            \item Создание простых и сложных таблиц
            \item Форматирование содержимого
            \item Отладка ошибок
            \item Профессиональное оформление документов
        \end{itemize}
    \end{enumerate}
\end{frame}

\begin{frame}{Практические рекомендации}
    \begin{block}{Для успешной работы с таблицами в LaTeX:}
        \begin{enumerate}
            \item Всегда проверяйте количество ячеек и столбцов
            \item Используйте подходящее выравнивание для типа данных
            \item Применяйте \texttt{multicolumn} для улучшения структуры
            \item Используйте линии для повышения читаемости
            \item Тестируйте таблицу после каждого изменения
        \end{enumerate}
    \end{block}
\end{frame}

\begin{frame}{Области применения}
    \begin{columns}[T]
        \begin{column}{0.5\textwidth}
            \begin{block}{Научные работы}
                \begin{itemize}
                    \item Статьи
                    \item Диссертации
                    \item Отчёты
                    \item Публикации
                \end{itemize}
            \end{block}
        \end{column}
        
        \begin{column}{0.5\textwidth}
            \begin{block}{Деловые документы}
                \begin{itemize}
                    \item Финансовые отчёты
                    \item Презентации
                    \item Аналитика
                    \item Статистика
                \end{itemize}
            \end{block}
        \end{column}
    \end{columns}
    
    \vspace{0.5cm}
    
    \begin{exampleblock}{Преимущество LaTeX}
        Высокое качество типографского оформления и автоматическое позиционирование таблиц
    \end{exampleblock}
\end{frame}

% Заключение
\section{Заключение}

\begin{frame}{Заключение}
    \begin{block}{Результаты работы}
        \begin{itemize}
            \item Изучены основные принципы создания таблиц в LaTeX
            \item Освоены команды выравнивания и форматирования
            \item Исследовано поведение при различных ошибках
            \item Изучена команда \texttt{\textbackslash multicolumn}
            \item Созданы таблицы различной сложности
        \end{itemize}
    \end{block}
    
    \vspace{0.5cm}
    
    \begin{exampleblock}{Практическая значимость}
        Полученные навыки позволяют создавать профессионально оформленные документы с таблицами любой сложности для научных и деловых целей
    \end{exampleblock}
\end{frame}

\begin{frame}{Список использованных источников}
    \begin{thebibliography}{99}
        \bibitem{lvovsky} Львовский С.М. Набор и вёрстка в системе LaTeX. — М.: МЦНМО, 2003.
        
        \bibitem{kotelnikov} Котельников И.А., Чеботаев П.З. LaTeX 2ε по-русски. — Новосибирск: Сибирский хронограф, 2004.
        
        \bibitem{latex} The LaTeX Project. Official Documentation. — \texttt{https://www.latex-project.org/}
        
        \bibitem{overleaf} Overleaf Documentation. Tables. — \texttt{https://www.overleaf.com/learn/latex/Tables}
        
        \bibitem{wikibooks} Wikibooks. LaTeX/Tables. — \texttt{https://en.wikibooks.org/wiki/LaTeX/Tables}
    \end{thebibliography}
\end{frame}

% Финальный слайд
\begin{frame}[plain]
    \begin{center}
        \Huge Спасибо за внимание!
    \end{center}
\end{frame}

\end{document}
